\section{Discussion}
\subsection{What the comparisons can establish}
The experimental structure was chosen to make a narrow but useful conclusion possible. If the morphology arm performs better than real-only training, the result supports an augmentation effect for the reviewed data, the chosen generator and the tested classifier configuration. It would not establish that explicit morphology caused the improvement unless the morphology arm also differs from the matched generic arm. Conversely, if generic and morphology images produce similar scores, increasing prompt detail may have no incremental value even when both improve over E0. These distinctions are important because generation quality, dataset diversity and model optimization are related but different quantities.

Quality and diversity selection answer another question. If E3 differs from E2 when both draw the same number of images from the same morphology pool, a plausible explanation is the selected samples rather than prompt formulation. That explanation still needs caution: a reviewer may favor familiar or visually conspicuous lesions, making a quality-ranked subset less representative of ordinary clinical presentations. Diversity-based selection can also overemphasize unusual but low-fidelity candidates. The selection policy must therefore be reported with the acceptance criteria and classwise sample counts, including cases where it fails to improve real-test performance.

\subsection{The synthetic-to-real gap}
A synthetic-only classifier may learn visually consistent lesion patterns yet perform poorly on real photographs. This outcome is not contradictory. Generated images can have cleaner backgrounds, similar face framing, recurring skin texture or category-specific color artifacts absent from the target collection. Combining real and synthetic data may reduce such mismatch, but it may also cause the network to rely on generator-specific shortcuts. A source-separated confusion analysis is therefore more informative than declaring a generator successful because its images look medically plausible.

A gain that appears only on photographs from one source is a warning sign. The study can describe that effect, but it should not call it evidence of broad clinical generalization. The collected real test is not equivalent to an independent prospective hospital cohort. A stronger external-validity claim would require a separately sourced, diagnosis-reviewed set acquired under a different imaging process and frozen before tuning.

\subsection{Threats to internal validity}
Several practical limitations are anticipated. First, the broad six-category taxonomy simplifies clinical distinctions. A premalignant lesion and an invasive carcinoma should not be conflated without explicit endpoint design. Second, labels imported from filenames or public galleries can be wrong or incomplete. The review gate reduces known errors but cannot replace qualified diagnostic assessment. Third, source and watermark artifacts may remain correlated with labels. Fourth, visual near-duplicate grouping cannot fully establish independence among patients, and incomplete grouping may leave correlated images across the split.

Fifth, base generation and LoRA adaptation are different interventions. If their outputs are mixed, any classifier difference may reflect the adapter, the text prompt or the sample-selection policy. The analysis must preserve the separate generator identity. Sixth, normal-skin synthesis uses SDXL while most disease generations use SD2.1; this architecture difference may introduce a category-correlated synthetic style that resembles a class shortcut. Finally, training-seed variation and uncertainty from a limited real test are different sources of variability. Repeating training three times does not create three independent test sets.

\subsection{Scope of clinical interpretation}
Even an improvement supported by paired real-test intervals would not justify using the classifier to diagnose an individual. Clinical use requires reliable diagnoses, clinically meaningful differential endpoints, diverse patient populations, error analysis with qualified experts and prospective testing under realistic acquisition conditions. Subgroup statements require dependable annotations and sufficient representation; appearance cannot be used to infer sensitive personal traits. The study is primarily a controlled analysis of synthetic augmentation under a specified photographic task.

\subsection{Discussion to finalize with measurements}
\noindent\Pending. Insert the exact observed outcome for each research question here. Explain whether E2--E1 isolates a consistent morphology-specific gain, whether quality/diversity selection helped at matched sample counts, which classes improved or deteriorated, and whether changes persist across architectures and seeds. Discuss observed domain shift and computational costs, including unsuccessful runs and negative findings. Update only after the statistics, source audit and clinical annotation quality are checked.
